%%
%% spring-lecture07.tex
%% 
%% Made by Alex Nelson <pqnelson@gmail.com>
%% Login   <alex@lisp>
%% 
%% Started on  2026-04-16T07:55:17-0700
%% Last update 2026-04-16T07:55:17-0700
%% 

\lecture[Levi-Civita Connection]{}

\begin{definition}
Let $(M, g)$ be a Riemannian manifold, let $\nabla$ be a connection on
$M$. We say the connection $\nabla$ is \define{Metric-Compatible} (or
``compatible with the metric $g$'') if for all vector fields
$X,Y,Z\in\Vect{M}$ we have $Xg(Y,Z)=g(\nabla_{X}Y,Z)+g(Y,\nabla_{X}Z)$
which should be viewed as a sort of Leibniz rule.
\end{definition}

\begin{theorem}[Fundamental theorem of Riemannian geometry]\label{thm:fundamental:riemannian-geometry}
Let $(M, g)$ be a Riemannian manifold. There exists a unique
torsion-free connection $\nabla$ on $M$ which is compatible with the
metric $g$.
\end{theorem}

\begin{proof}
\textsc{Claim 1: Uniqueness.} Assume we have such a connection
$\nabla$. Let $X,Y,Z\in\Vect{M}$ be arbitrary vector fields on $M$. We
have (by metric compatibility):
\begin{subequations}
  \begin{align}
X(g(Y,Z)) &= g(\nabla_{X}Y,Z) + g(Y,\nabla_{X}Z)\\ 
Y(g(Z,X)) &= g(\nabla_{Y}Z,X) + g(Z,\nabla_{Y}X)\\ 
Z(g(X,Y)) &=g(\nabla_{Z}X,Y) + g(X,\nabla_{Z}Y), 
  \end{align}
\end{subequations}
Then we have (underlining terms we want to eventually cancel out):
\begin{multline}\label{eq:fundamental-thm-riemannian-geometry:step-1}
X(g(Y,Z))+Y(g(Z,X))-Z(g(X,Y))\\
=g(\nabla_{X}Y,Z) + \underline{g(Y,\nabla_{X}Z)} + \underline{g(\nabla_{Y}Z,X)} + g(Z,\nabla_{Y}X)
- g(\nabla_{Z}X,Y) - g(X,\nabla_{Z}Y)
\end{multline}
Using torsion-free, we have
\begin{subequations}
  \begin{align}
\nabla_{Y}X &= \nabla_{X}Y + [Y,X]\\
\nabla_{Z}X &= \nabla_{X}Z + [Z,X]\\
\nabla_{Z}Y &= \nabla_{Y}Z + [Z,Y]
  \end{align}
\end{subequations}
which we use to rewrite Equations~\eqref{eq:fundamental-thm-riemannian-geometry:step-1}
leaving only the first term untouched:

\begin{subequations}
  \begin{align}
&X(g(Y,Z))+Y(g(Z,X))-Z(g(X,Y))\nonumber\\
&=g(\nabla_{X}Y,Z) + \underline{g(Y,\nabla_{X}Z)} + \underline{g(\nabla_{Y}Z,X)} + g(Z,\nabla_{X}Y + [Y,X])\nonumber\\
&\phantom{=\quad}- g(\nabla_{X}Z + [Z,X],Y) - g(X,\nabla_{Y}Z + [Z,Y])\\
&=g(\nabla_{X}Y,Z) + \underline{g(Y,\nabla_{X}Z)} + \underline{g(\nabla_{Y}Z,X)} + g(Z,\nabla_{X}Y) + g(Z,[Y,X])\nonumber\\
&\phantom{=\quad}- g(\nabla_{X}Z, Y) - g([Z,X],Y) - g(X,\nabla_{Y}Z) - g(X, [Z,Y])\\
&=2g(\nabla_{X}Y,Z) + g(Z,[Y,X]) - g([Z,X],Y)  - g(X, [Z,Y]),
  \end{align}
\end{subequations}
where we use the fact the metric is symmetric in its arguments and linear.
Then we obtain
\begin{equation}\label{eq:fundamental-thm-riemannian-geometry:key-result}
\begin{split}
g(\nabla_{X}Y,Z) &= \frac{1}{2}\biggl(X(g(Y,Z))+Y(g(Z,X))-Z(g(X,Y))\\
&\phantom{=\frac{1}{2}\biggl(}- g(Z,[Y,X]) + g([Z,X],Y)  + g(X, [Z,Y])\biggr)
\end{split}
\end{equation}
Therefore $\nabla_{X}Y$ is uniquely determined by the metric $g$.

\textsc{Claim 2: Existence.} It is sufficient to check the right-hand
side of
Equation~\eqref{eq:fundamental-thm-riemannian-geometry:key-result} is
$C^{\infty}(M)$-linear in $Z$, and using
Equation~\eqref{eq:fundamental-thm-riemannian-geometry:key-result} as
a definition we obtain a torsin-free connection which is compatible
with $g$.
\end{proof}

\begin{definition}
Let $(M,g)$ be a Riemannian manifold.
The connection $\nabla$ given by the previous theorem (i.e., the
unique torsion-free connection compatible with the metric) is called
the \define{Levi-Civita Connection}.
\end{definition}

\begin{definition}
Let $(M,g)$ be a Riemannian $n$-manifold. Let $\nabla$ be a connection
on $M$. Let $(U,\varphi=(\xi^{1},\dots,\xi^{n}))\in\smoothstr{M}$ be a
chart in the smooth structure of $M$. We define the
\define{Christoffel Symbols} of $\nabla$ in the local chart
$(U,\varphi)$ to be the smooth functions $\Gamma^{k}_{ij}\in C^{\infty}(U)$
satisfying
\begin{equation}
\nabla_{\partial/\partial\xi^{i}}\frac{\partial}{\partial\xi^{j}}=\sum^{n}_{k=1}\Gamma^{k}_{ij}\frac{\partial}{\partial\xi^{k}}.
\end{equation}
\end{definition}

\begin{node}
Christoffel symbols are a reasonable thing to look for: covariant
differentiation of a vector field \emph{is} a vector field; and if we
have a basis for vector fields (the coordinate vector fields), then it is
reasonable to check that the covariant derivative of the coordinate
vector field with respect to itself results in a linear combination of
the coordinate vector fields. The coefficients of this linear
combination are precisely the Christoffel symbols.
\end{node}

\begin{remark}
We want, in our chart $(U,\varphi)$, for every vector fields $X,Y\in\Vect{M}$
where
\begin{equation}
X=\sum^{n}_{i=1}X^{i}\frac{\partial}{\partial\xi^{i}},\quad\mbox{and}\quad
Y=\sum^{n}_{j=1}Y^{j}\frac{\partial}{\partial\xi^{j}}.
\end{equation}
Then
\begin{subequations}
  \begin{align}
\nabla_{X}Y &= \nabla_{\sum_{i}X^{i}\partial/\partial\xi^{i}}\sum^{n}_{j=1}Y^{j}\frac{\partial}{\partial\xi^{j}}\\
&=\sum^{n}_{i=1}X^{i}\nabla_{\partial/\partial\xi^{i}}\left(\sum^{n}_{j=1}Y^{j}\frac{\partial}{\partial\xi^{j}}\right)\quad\mbox{by $C^{\infty}$-linearity of first argument of $\nabla$}\\
&=\sum^{n}_{i,j=1}X^{i}\nabla_{\partial/\partial\xi^{i}}\left(Y^{j}\frac{\partial}{\partial\xi^{j}}\right)
\quad\mbox{by additivity of second argument of $\nabla$}\\
&=\sum_{i,j}X^{i}\frac{\partial Y^{j}}{\partial\xi^{i}}\frac{\partial}{\partial\xi^{j}}+X^{i}Y^{j}\nabla_{\partial/\partial\xi^{i}}\frac{\partial}{\partial\xi^{j}}
\quad\mbox{by Leibniz rule}\\
&=\sum_{i,j}\left(X^{i}\frac{\partial Y^{j}}{\partial\xi^{i}}\frac{\partial}{\partial\xi^{j}}+\sum_{k}X^{i}Y^{j}\Gamma^{k}_{ij}\frac{\partial}{\partial\xi^{k}}\right)\quad\mbox{by def.\ of Christoffel symbols}\\
&=\sum_{k=1}^{n}\left(\sum^{n}_{i}X^{i}\frac{\partial Y^{k}}{\partial\xi^{i}}+\sum^{n}_{i,j=1}\Gamma^{k}_{ij}X^{i}Y^{j}\right)\frac{\partial}{\partial\xi^{k}}
  \end{align}
\end{subequations}
where we reindex to get the last line.

Observe, we are just taking the first derivatives of $Y$, and in this
sense the connection $\nabla$ is a first-order differential
operator. The first term $X^{i}\partial Y^{k}/\partial\xi^{i}$ is the
``Euclidean'' or ``flat part''. Compare it to the Euclidean
connection's $\sum_{k}\langle\nabla Y^{k},X\rangle\partial/\partial\xi^{k}$.
The Christoffel symbols are a perturbation of the intuitive thing,
which is bilinear in both $X$ and $Y$.
\end{remark}

\begin{node}
On a submanifold of $\RR^{k}$, what happens? There is an intuitive way
to get a connection. Recall the tangent space for a smooth submanifold
$M$ of $\RR^{k}$ will be canonically identified with a linear subspace
of $\RR^{k}$. If you differentiate a tangent vector field on $M$ in
this way, you could ``leave'' the tangent space.

There is a solution: take the projection (of the result) back onto the
tangent space. That is \emph{precisely} what the Levi-Civita term is
doing! Let us make this intuition more precise.
\end{node}

\begin{example}
Let $M\subset\RR^{k}$ be a smooth submanifold. Then we can define
\begin{equation}\label{eq:spring-lecture07:math151c:induced-connection-from-nash-embedding}
(\nabla_{X}Y)(x)=\pi_{x}\bigl((\nabla^{0}_{X}Y)(x)\bigr)
\end{equation}
where $\pi_{x}\colon\RR^{k}\to T_{x}$ is the projection done pointwise
for $x\in M$, and $\nabla^{0}$ is the Euclidean connection on
$\RR^{k}$. Is the Levi-Civita connection for the induced metric on $M$
the same thing as the induced connection from this equation?
\end{example}

\begin{xca}
Check this is true! (Check Equation~\eqref{eq:spring-lecture07:math151c:induced-connection-from-nash-embedding}
is torsion-free and compatible with the metric on $M$ induced by $\RR^{k}$)
\end{xca}

\begin{node}
We can extend covariant differentiation of \emph{vector fields} to
covariant differentiation of \emph{tensors}. Recall tensors are
sections of tensor bundles. A first step is to ask for covariant
differentiation of 1-forms, then extend by Leibniz rule for tensor products.
\end{node}

\begin{definition}[Extension of $\nabla$ to $T^{*}M$]
Let $\nabla$ be a connection on a smooth manifold $M$, let
$\omega\in\Omega^{1}(M)$ be a covector field (``one-form''). Then for
every $X\in\Vect{M}$, we define the \define{Covariant Derivative} of
$\omega$ with respect to $X$, denoted $\nabla_{X}\omega$, to be the
one-form such that for every $Y\in\Vect{M}$ we have
\begin{equation}
(\nabla_{X}\omega)[Y]:=\nabla_{X}(\omega[Y])-\omega(\nabla_{X}Y),
\end{equation}
but since $\omega[Y]\in C^{\infty}(M)$, this is really equivalent to
\begin{equation}
(\nabla_{X}\omega)[Y]:=X(\omega[Y])-\omega(\nabla_{X}Y).
\end{equation}
This defines a connection on $T^{*}M$. We can check: 
\begin{enumerate}
\item $\nabla_{X}\omega$ is $C^{\infty}(M)$-linear in $X$:
  $\nabla_{fX}\omega=f\nabla_{X}\omega$ for any $f\in C^{\infty}(M)$; and
\item $\nabla_{X}\omega$ is $\RR$-linear in $\omega$:
  $\nabla_{X}(a\phi+b\omega)=(a\nabla_{X}\phi)+(b\nabla_{X}\omega)$
  for any $a,b\in\RR$ and $\phi,\omega\in T^{*}M$; and
\item Leibniz rule: $\nabla_{X}(f\omega)=X(f)\omega+f\nabla_{X}\omega$
  for any $f\in C^{\infty}(M)$
\end{enumerate}
\end{definition}